\newpage
\phantomsection
\label{prf:induction-principle-for-peano-system}

\begin{remark*}[Return]
\hyperref[thm:induction-principle-for-peano-system]{Return to Theorem}
\end{remark*}

\begin{theorem*}[Induction Principle for a Peano System]
Let $(P,S,1)$ be a Peano system, and let $\varphi$ be a predicate on $P$.
If
\[
\varphi(1)
\qquad\text{and}\qquad
\forall n \in P\;(\varphi(n) \Longrightarrow \varphi(S(n))),
\]
then
\[
\forall n \in P\;\varphi(n).
\]
\end{theorem*}

\begin{proof}
\textbf{Professional Standard Proof.}~\\
Let
\[
A:=\{n \in P : \varphi(n)\}.
\]
We will verify that $A$ satisfies the hypotheses of the induction axiom
for the Peano system $(P,S,1)$.

Because $\varphi(1)$ holds, we have
\[
1 \in A.
\]

Now let $n \in P$ and suppose that $n \in A$. By the definition of $A$,
this means that $\varphi(n)$ holds. By the inductive hypothesis
\[
\forall n \in P\;(\varphi(n) \Longrightarrow \varphi(S(n))),
\]
it follows that $\varphi(S(n))$ holds. Therefore
\[
S(n) \in A.
\]
So
\[
\forall n \in P\;(n \in A \Longrightarrow S(n) \in A).
\]

Since $(P,S,1)$ is a Peano system, Axiom~\ref*{ax:peano-induction} applies.
Hence
\[
A=P.
\]

Finally, let $n \in P$. Because $A=P$, we have $n \in A$. By the definition
of $A$, this means that $\varphi(n)$ holds. Since $n \in P$ was arbitrary,
we conclude that
\[
\forall n \in P\;\varphi(n).
\]
This is exactly the desired conclusion.
\end{proof}

\begin{proof}
\textbf{Detailed Learning Proof.}~\\

\textbf{Step 1.} Build the subset that corresponds to the predicate.

The induction axiom in the definition of a Peano system is stated for subsets
of the underlying set $P$, while the theorem we want is stated for a
predicate $\varphi$ on $P$. To connect these two forms, define
\[
A:=\{n \in P : \varphi(n)\}.
\]
So $A$ is exactly the set of elements of $P$ for which the property
$\varphi$ holds.

\textbf{Step 2.} Check the base case in subset language.

We are given that
\[
\varphi(1)
\]
holds. By the definition of $A$, that means precisely that
\[
1 \in A.
\]
So the first hypothesis required by the induction axiom is satisfied.

\textbf{Step 3.} Check closure under successor in subset language.

Now let $n \in P$ and assume that
\[
n \in A.
\]
By the definition of $A$, this means that
\[
\varphi(n)
\]
holds. But the theorem also assumes that
\[
\forall n \in P\;(\varphi(n) \Longrightarrow \varphi(S(n))).
\]
Therefore $\varphi(S(n))$ holds, and so again by the definition of $A$ we get
\[
S(n) \in A.
\]
Since $n \in P$ was arbitrary, we have shown that
\[
\forall n \in P\;(n \in A \Longrightarrow S(n) \in A).
\]
Thus $A$ is closed under the successor operation.

\textbf{Step 4.} Apply the induction axiom of the Peano system.

We have now proved both hypotheses needed for
\hyperref[ax:peano-induction]{Axiom~\ref*{ax:peano-induction}}:
\[
1 \in A
\qquad\text{and}\qquad
\forall n \in P\;(n \in A \Longrightarrow S(n) \in A).
\]
Because $(P,S,1)$ is a Peano system, the induction axiom tells us that
\[
A=P.
\]

\textbf{Step 5.} Translate back to the predicate form.

To prove the desired conclusion, let $n \in P$. Since $A=P$, we must have
\[
n \in A.
\]
By the definition of $A$, this means that $\varphi(n)$ holds.

Because $n \in P$ was arbitrary, we conclude that
\[
\forall n \in P\;\varphi(n).
\]

\textbf{Conclusion.} The subset-form induction axiom implies the usual
predicate-form induction principle.
\end{proof}

\begin{remark*}[Proof structure]
The proof converts the predicate $\varphi$ into the subset
\[
A=\{n \in P : \varphi(n)\}.
\]
The base case $\varphi(1)$ becomes the statement $1 \in A$, and the
successor step for $\varphi$ becomes closure of $A$ under the successor map.
Once these two facts are established, the induction axiom for the Peano
system forces $A=P$. Translating back from membership in $A$ to truth of
$\varphi$ yields the predicate-form induction principle.
\end{remark*}

\begin{remark*}[Dependencies]~\\
\begin{itemize}
  \item \hyperref[def:peano-system]{Definition~\ref*{def:peano-system}}
  \item \hyperref[ax:peano-induction]{Axiom~\ref*{ax:peano-induction}}
\end{itemize}
\end{remark*}
